% Generated draft: requires mathematical and source review.
\documentclass[11pt]{article}
\usepackage[margin=25mm]{geometry}
\usepackage{fontspec}
\setmainfont{Latin Modern Roman}
\usepackage{amsmath,amssymb,mathtools,unicode-math}
\setmathfont{Latin Modern Math}
\usepackage{enumitem,needspace,xurl,hyperref}
\setlength{\parindent}{0pt}
\setlength{\parskip}{5pt}
\setlength{\emergencystretch}{3em}
\newcommand{\N}{\mathbb{N}}
\newcommand{\Z}{\mathbb{Z}}
\newcommand{\Q}{\mathbb{Q}}
\newcommand{\R}{\mathbb{R}}
\newcommand{\C}{\mathbb{C}}
\newcommand{\F}{\mathbb{F}}
\newcommand{\A}{\mathbb{A}}
\newcommand{\PP}{\mathbb{P}}
\newcommand{\E}{\mathbb{E}}
\newcommand{\eps}{\varepsilon}
\newcommand{\dd}{\,\mathrm{d}}
\begin{document}
\section*{UnsolvedMath ID: 800582}
\textbf{Source problem code:} catalogue entry
\subsection{Definitions and mathematical statement}
The Continuum Hypothesis is the assertion
\[
\mathrm{CH}:\quad 2^{\aleph_0}=\aleph_1.
\]
For families \(\mathcal A\) and \(\mathcal D\) of collections of subsets of a set \(X\), the selection principle \(\mathsf S_1(\mathcal A,\mathcal D)\) means that for every sequence \((\mathcal U_n)_{n\in\N}\) of members of \(\mathcal A\), there are choices \(U_n\in\mathcal U_n\) such that \(\{U_n:n\in\N\}\in\mathcal D\). The principle \(\mathsf S_{\mathrm{fin}}(\mathcal A,\mathcal D)\) means that for every such sequence there are finite subfamilies \(\mathcal F_n\subseteq\mathcal U_n\) such that \(\bigcup_{n\in\N}\mathcal F_n\in\mathcal D\). These abstract definitions do not specify the particular families denoted by \(\mathrm B\), \(\mathrm B_\Gamma\), and \(\mathrm B_\Omega\) in the source; their definitions and cover conventions are absent from the supplied record.

The question attributed to Tsaban (2006), Problem 8, is whether \(\mathrm{CH}\) implies that each of Settings \((d),(e),(g),(j),(l)\) can be realized. Its logical form can only be recorded schematically as
\[
\mathrm{CH}\ \Longrightarrow\
\bigwedge_{s\in\{d,e,g,j,l\}}\exists X_s\;\operatorname{Realizes}(X_s,s).
\]
Here \(\operatorname{Realizes}\) is an unresolved placeholder: the supplied extraction does not define any of the named settings or the domain of their realizations. Thus this display records the requested separate realizations, but is not a complete mathematical formulation. In particular, it does not require one object to realize all five settings simultaneously.

The source supplies the following context under \(\mathrm{CH}\). It reports realizations of Settings \((c),(h),(i)\) in [coc2], of Setting \((k)\) in [o-bdd], and of Setting \((n)\) in [Rec94, Brendle96, MilNonGamma]. For Setting \((b)\), it proposes taking \(L\) as in Setting \((c)\), \(S\) as in Setting \((i)\), and \(X=L\cup S\). It states that additivity of \(\mathsf S_{\mathrm{fin}}(\mathrm B,\mathrm B)\) gives this property for \(X\), whereas heredity for subsets of \(\mathsf S_1(\mathrm B_\Gamma,\mathrm B_\Omega)\) and \(\mathsf S_1(\mathrm B,\mathrm B)\), cited to [ideals], shows that \(X\) satisfies neither of these latter properties. The missing definitions of Settings \((c)\) and \((i)\) prevent reconstruction of the hypotheses on \(L\) and \(S\). Finally, the source tentatively suggests that forcing arguments similar to those of [Brendle96] could realize Settings \((f)\) and \((m)\); this suggestion is not presented as an established realization.
\subsection{Short English statement}
Assuming the Continuum Hypothesis, must each of the five configurations labeled \((d),(e),(g),(j),(l)\) have a realization? The supplied record omits the definitions of these configurations and the permitted objects, so the intended problem cannot yet be stated completely. The registry labels it open; this supplied status has not been independently verified.
\subsection{Sources}
\begin{itemize}
\item Local UnsolvedMath registry: permanent ID \(800582\), exact source code \texttt{catalogue entry}, title ``Selection Principles in Mathematics --- Problem 8,'' in the AMR Open Problem Lists collection. The record attributes the problem to Tsaban, with proposed year 2006, and supplies status ``open.'' Its creation and update timestamps are both \texttt{2026-10-16T14:00:00Z}; upstream revision: \texttt{d371ee917ac9647fbb06cd6fc529deec6f3dda82}.
\item Supplied source attribution: Tsaban, \emph{Selection Principles and Special Sets of Reals: Open Problems} (2006), Problem 8. Supplied URL: \url{https://arxiv.org/abs/math/0609601}. The provenance records access on 2026-10-16 and extraction method \texttt{arxiv\_latex}. These are supplied provenance details, not a claim of independent access or verification.
\item Unresolved reference labels retained from the supplied excerpt: [coc2], [o-bdd], [Rec94], [Brendle96], [MilNonGamma], and [ideals]. The record supplies no bibliographic entries or links for these labels. Its realization claims, hereditary-property attribution, and tentative forcing suggestion are preserved as source reports.
\end{itemize}
\end{document}
